\chapter{Sal de la tierra}

\epigraph{«Vosotros sois la sal de la tierra. [\ldots] Vosotros sois la luz del mundo.»}{Mateo 5,13-14}

Si hay algo que la gente cree saber sobre lo que la Iglesia piensa del matrimonio, es esto: que quiere que tengas muchos hijos. Cuantos más, mejor. Sin plantearse nada, sin pensárselo dos veces. Abrir la puerta a la vida ---dicen--- significa no cerrarla nunca, como si la fecundidad fuera una cuestión de aritmética y el valor de un matrimonio se midiera por el número de sillas que necesita alrededor de la mesa.

La caricatura es conocida. Y como toda caricatura, tiene un rasgo reconocible ---algo de verdad--- rodeado de exageración. Sí, la Iglesia valora la apertura a la vida. Sí, los hijos son un bien enorme. Pero reducir la fecundidad del matrimonio a la cantidad de hijos es como reducir la calidad de un árbol al número de frutos que da en una temporada, sin mirar si da sombra, si sus raíces sujetan la tierra, si sus ramas acogen nidos.

La fecundidad de un matrimonio es algo mucho más ancho, mucho más hondo y mucho más hermoso de lo que esa caricatura sugiere. Y para verlo, hay que cambiar la mirada.

\section{La familia que todos recuerdan}

Dejadme que os cuente algo que probablemente habéis vivido. Pensad un momento en las familias que habéis conocido a lo largo de vuestra vida. No las familias perfectas ---esas no existen fuera de los anuncios de Navidad---. Las familias \emph{reales}. Las que os han marcado.

Seguro que hay alguna. Una casa a la que siempre os apetecía ir. Donde os sentíais acogidos sin necesidad de avisar. Donde la puerta estaba abierta y la mesa se estiraba sin drama para un comensal más. Una casa donde se reían mucho, donde los padres discutían a veces pero se miraban con una complicidad que daba gusto ver. Una familia donde los amigos de los hijos acababan quedándose a dormir y, con el tiempo, siendo casi de la familia.

¿La tenéis? Bien. Ahora fijaos en algo curioso: esa familia que recordáis con cariño, esa familia que os hizo sentir que el mundo era un lugar más habitable, probablemente no hacía nada espectacular. No organizaba grandes eventos. No tenía un programa de acción social. Simplemente \emph{vivía} de una manera que generaba vida a su alrededor. Como la sal, que no se ve en el plato pero cambia todo el sabor.

Eso es fecundidad.

\section{¿Y si fuera al revés?}

Porque la percepción habitual ---fecundidad igual a hijos, punto--- no solo es reduccionista. Es que, a veces, el enfoque se invierte hasta decir lo contrario de lo que debería.

Pensemos en un matrimonio que tiene cinco hijos pero que vive encerrado en sí mismo. Una familia que no se abre, que no acoge, que no comparte. Una casa donde no entra nadie de fuera. Unos padres que no tienen tiempo para nadie más ---ni para ellos mismos---, no porque estén ocupados dándose, sino porque han convertido la familia en una fortaleza con el puente levadizo siempre alzado. ¿Es fecundo ese matrimonio? Tiene frutos, sí. Pero ¿da vida?

Ahora pensemos en otro matrimonio. Quizá tiene dos hijos, quizá tiene uno, quizá ---por razones que solo ellos y Dios conocen--- no ha podido tener ninguno. Pero es una casa donde la gente quiere estar. Una pareja que acompaña a otros matrimonios en crisis. Unos esposos que acogen al sobrino que lo está pasando mal, que visitan al vecino enfermo, que abren su salón para que un grupo de amigos pueda reunirse a rezar o simplemente a compartir la vida. Una familia que \emph{irradia} algo que los demás reconocen sin poder explicarlo del todo: alegría, sentido, calor, esperanza.

¿Cuál de los dos es más fecundo?

La pregunta no pretende enfrentar a las familias grandes con las pequeñas. No se trata de eso. Se trata de entender que la fecundidad no se mide en números, sino en vida generada. En cuánta vida brota alrededor de ese matrimonio. En cuántas personas encuentran en esa familia un refugio, un referente, una razón para creer que el amor de verdad es posible.\footnote{«La fecundidad del amor matrimonial no se reduce a la sola procreación de los hijos, sino que debe extenderse a su educación moral y a su formación espiritual. [\ldots] Los esposos, a los que Dios no ha concedido tener hijos, pueden, no obstante, tener una vida conyugal plena de sentido, humana y cristianamente.» \textit{Catecismo de la Iglesia Católica}, n. 1654.}

\section{El amor que desborda}

Hay una ley no escrita que se cumple siempre: el amor verdadero no cabe dentro.

Cuando dos personas se aman de verdad ---no con el amor de las películas, que mira al otro como si el resto del mundo no existiera, sino con el amor maduro, el que elige y permanece---, ese amor tiende a salir. A desbordarse. A buscar a quién más darse. No porque los esposos se aburran el uno del otro, sino porque el amor es así: cuanto más auténtico, más expansivo. Cuanto más hondo, más ancho.

Los hijos son la expresión más natural de ese desbordamiento. No nacen porque «hay que tenerlos», ni porque «la Iglesia manda». Nacen porque el amor de dos personas se hace tan real, tan encarnado, que se convierte en una tercera persona. El hijo no es un proyecto ni una obligación: es el amor que cobra rostro, el abrazo que se hace carne.\footnote{«El hijo no viene de fuera a añadirse al amor mutuo de los esposos; brota del corazón mismo de ese don recíproco, del que es fruto y cumplimiento.» Juan Pablo II, Exhortación apostólica \textit{Familiaris Consortio}, n. 14 (1981).} Hay algo sobrecogedor en eso, si se piensa despacio: que dos personas puedan amarse tanto que de su amor nazca alguien nuevo, alguien que nunca antes existió, alguien que vivirá para siempre. Pocas cosas en la experiencia humana se acercan más al acto creador de Dios.

Pero el desbordamiento no se agota en los hijos. El amor fecundo sigue buscando cauces, sigue inventando formas de dar vida.

\section{La fecundidad que no se cuenta}

La tradición cristiana ha hablado siempre de varias dimensiones de la fecundidad, aunque a veces las hayamos olvidado.\footnote{El Concilio Vaticano II amplía la comprensión de la fecundidad matrimonial más allá de la procreación: «El matrimonio y el amor conyugal están ordenados por su propia naturaleza a la procreación y educación de la prole. [\ldots] Sin embargo, el matrimonio no ha sido instituido solamente para la procreación.» \textit{Gaudium et Spes}, n. 48 y 50.}

Hay una \textbf{fecundidad educativa}: la que forma personas. No basta con traer hijos al mundo; hay que ayudarles a crecer, a pensar, a ser libres, a amar. Un padre que enseña a su hijo a pedir perdón está siendo fecundo. Una madre que le enseña a mirar al que sufre está generando vida. Y no solo para sus hijos: la educación que una familia transmite se extiende a quienes la rodean, porque los valores se contagian como el acento.

Hay una \textbf{fecundidad espiritual}. Es la del matrimonio que reza junto, que busca a Dios, que transmite la fe no como una lección teórica sino como algo vivo, experimentado, agradecido. Es también la del matrimonio que acompaña a otros en su camino de fe ---parejas que preparan a novios para el matrimonio, que animan grupos de oración, que simplemente testimonian con su vida que creer tiene sentido---.\footnote{El papa Francisco habla de la «fecundidad ampliada» de la familia: «Una familia que no acoge y no sale de sí misma deja de ser familia.» Cf. Francisco, Exhortación apostólica \textit{Amoris Laetitia}, n. 181 (2016).}

Hay una \textbf{fecundidad social}: la familia que se implica en su comunidad, que participa, que aporta. La que acoge al inmigrante, la que cuida al anciano, la que se compromete con lo que pasa en su barrio. La familia no es una isla. Es ---debería ser--- un foco de calor que derrite un poco el hielo de alrededor.

Y hay una fecundidad que quizá no tiene nombre técnico pero que es la más visible de todas: la \textbf{fecundidad del testimonio}. Es la que se produce cuando alguien mira a una pareja y piensa: «Yo quiero eso.» No porque esa pareja sea perfecta, sino porque es real. Porque se nota que se quieren, que se respetan, que luchan juntos, que se perdonan, que disfrutan estando juntos. Esa envidia sana que despiertan ciertos matrimonios es, quizá, la forma más poderosa de evangelización que existe. No hace falta dar discursos. No hace falta repartir folletos. Basta con \emph{vivir} de manera que la gente mire y vea que el amor para siempre no es una fantasía de ingenuos, sino algo posible, real, deseable.

\section{La casa de la puerta abierta}

Me gusta pensar en la familia fecunda como una casa con la puerta abierta. No abierta de par en par a cualquier hora, que eso no es acogida, es caos. Pero sí abierta de verdad: una casa donde el que llega sabe que es bienvenido.

Todos hemos conocido casas así. Casas donde el sofá siempre tiene sitio para uno más. Donde la conversación fluye y nadie mira el reloj. Donde los hijos traen amigos con naturalidad y los amigos traen a los amigos de los amigos. Casas donde se ha llorado y se ha reído, donde se ha discutido a gritos y se ha hecho las paces en voz baja, donde la vida ha entrado con todo su peso y toda su belleza.

Esas casas no nacen por casualidad. Nacen de matrimonios que han entendido ---a veces sin saberlo formular--- que el amor no es un tesoro para guardar en una caja fuerte, sino un fuego que solo se mantiene vivo si se comparte. Que cuanto más das, más tienes. Que la generosidad no empobrece: multiplica.

En el capítulo primero de este libro hablamos de algo que tal vez recordéis: la supervivencia de la especie humana depende de la cooperación.\footnote{Cf. capítulo 1 de este libro.} Nuestros bebés nacen tan inmaduros, tan vulnerables, que sin una red de cuidado no sobreviven. Y esa red ---madre, padre, abuelos, comunidad--- es la primera forma de familia. La biología nos dice que no estamos hechos para criar solos, para vivir solos, para amar solos.

Pues bien: la fecundidad del matrimonio es la contribución de esa familia a la red. Es el hilo que los esposos tejen en la trama grande de la comunidad humana. Cada familia fecunda ---fecunda de verdad, no solo en número--- fortalece el tejido entero. Cada matrimonio que irradia vida hace que el mundo sea un poco más habitable para todos.

\section{La generosidad que piensa}

Y aquí hay que hablar de algo que a veces se omite por miedo a decir lo incorrecto: la paternidad responsable.

Porque si la fecundidad no es solo cuestión de números, tampoco la generosidad es cuestión de imprudencia. La Iglesia no pide a los matrimonios que tengan «todos los hijos que vengan» sin más, como quien deja un grifo abierto y se desentiende. Lo que pide es algo más exigente y más libre a la vez: que sean generosos \emph{y} prudentes. Que acojan a los hijos como un don, sí, pero que lo hagan con responsabilidad, discerniendo juntos ---los dos, en diálogo, ante Dios--- cuántos hijos pueden amar bien.\footnote{«Por razones justificadas, los esposos pueden querer espaciar los nacimientos de sus hijos. [\ldots] La paternidad responsable comporta [\ldots] el dominio que la razón y la voluntad deben ejercer.» Pablo VI, Encíclica \textit{Humanae Vitae}, n. 10 y 16 (1968). Cf. también \textit{Gaudium et Spes}, n. 50: «Los esposos [\ldots] de común acuerdo y esfuerzo, se formarán un juicio recto, atendiendo tanto a su propio bien como al de los hijos, ya nacidos o todavía por venir.»}

No es lo mismo tener tres hijos con el corazón abierto, bien educados, bien queridos, que tener ocho por inercia, saturados de cansancio y sin espacio emocional para ninguno. Y no es lo mismo tener un hijo único al que se ama con generosidad y al que se enseña a compartir, que no tener ninguno por puro egoísmo, porque «los niños molestan» o porque «hay que vivir la vida».

La clave, como siempre en el matrimonio cristiano, no está en el número. Está en el corazón. ¿Está abierto o cerrado? ¿Da vida o la retiene? ¿Es generoso o calculador?

Un matrimonio que tiene dos hijos y vive abierto al mundo es más fecundo que uno que tiene diez y vive replegado sobre sí mismo. Un matrimonio sin hijos que adopta, que acoge, que derrama vida a su alrededor, puede ser extraordinariamente fecundo. La fecundidad no es un dato demográfico. Es una actitud del alma. Eso sí, que la fecundidad de la vida sea amplia no borra el significado propio de cada gesto de amor entre los esposos; son dos niveles distintos, y en el próximo capítulo volveremos sobre ello.

Y esto merece un momento de atención especial, porque a veces se pasa de largo demasiado rápido. Un matrimonio que no tiene hijos ---por las razones que sean, que pertenecen al misterio de cada vida--- no es un matrimonio con la fecundidad cerrada. Es, si cabe, un matrimonio con la fecundidad \emph{especialmente} abierta. Porque toda la capacidad de entrega, toda la energía de amor que esos esposos llevan dentro, no se queda sin cauce: busca otros caminos. Y los encuentra. La acogida del que sufre. El servicio al que necesita. La entrega generosa al amigo, al vecino, al desconocido. La irradiación callada de una vida compartida que, precisamente porque no se cierra sobre los hijos propios, se abre de par en par al mundo.

No tener hijos biológicos no clausura la fecundidad. Puede, paradójicamente, lanzarla de una manera que pocos imaginan. Hay matrimonios sin hijos que han sido padres y madres de medio barrio, que han sostenido comunidades enteras, que han dado más vida a su alrededor que muchas familias numerosas. Porque la fecundidad, como venimos diciendo, no se mide por la biología. Se mide por la vida que brota alrededor.

\section{Sal y luz}

Jesús usó dos imágenes muy concretas para hablar de lo que sus discípulos estaban llamados a ser: sal y luz.\footnote{Mt 5,13-14.}

La sal no se ve. Nadie dice, al probar un plato bien sazonado: «Qué buena está la sal.» Dice: «Qué bueno está el plato.» La sal desaparece para que el sabor aparezca. No se impone, no grita, no llama la atención. Simplemente \emph{está}, y su presencia cambia todo lo que toca.

La luz tampoco pide permiso. No necesita explicar lo que hace. Simplemente ilumina. Entra en la habitación oscura y la oscuridad retrocede, no porque la luz luche contra ella, sino porque la oscuridad no puede sostenerse ante la luz.

Un matrimonio fecundo es sal y luz. No necesita dar lecciones sobre el amor. No necesita demostrar nada. Solo necesita vivir lo que es: una alianza de amor que genera vida. Y cuando lo hace, la gente lo nota. Los amigos lo notan. Los vecinos lo notan. Los hijos ---propios y ajenos--- lo notan. Y algo se enciende en ellos: el deseo de vivir así. La intuición de que eso es posible. La esperanza de que el amor no es una mentira.

Esa es la fecundidad más profunda del matrimonio. No la que se ve en el certificado de familia numerosa. La que se ve en los ojos de quienes han sido tocados por esa familia. La que germina en silencio, como la semilla en la tierra, y da fruto treinta, sesenta, cien por uno.\footnote{Cf. Mt 13,8.}

\section{Lo que el mundo necesita ver}

Vivimos en un tiempo raro. Un tiempo en el que el amor se consume rápido y se descarta fácil. En el que las relaciones se miden por lo que me aportan y se abandonan cuando dejan de ser rentables. En el que la palabra «para siempre» se ha vuelto sospechosa, como una promesa que nadie se atreve a tomarse en serio.

En ese contexto, un matrimonio que dura, que crece, que se abre a los demás, que genera vida en todas sus formas, es un acto casi revolucionario. No porque sea perfecto ---ya lo hemos dicho: la perfección no existe en la vida real---. Sino porque es \emph{posible}. Porque demuestra, con hechos más que con palabras, que el amor fiel, indisoluble y fecundo no es una reliquia del pasado ni una imposición de la Iglesia: es lo que el corazón humano desea en lo más hondo.

Cada matrimonio que vive así ---con sus luchas, con sus caídas, con sus reconciliaciones--- está diciendo al mundo algo que el mundo necesita oír desesperadamente: que vale la pena apostar por el amor. Que vale la pena quedarse. Que vale la pena dar vida.

No hace falta ser héroes. No hace falta ser santos de estampa. Basta con ser sal. Basta con ser luz. Basta con vivir el amor con el corazón abierto y dejar que ese amor haga su trabajo.

Porque un matrimonio fecundo no es el que llena la casa de hijos. Es el que llena el mundo de esperanza.

\paracaminar

\begin{enumerate}
  \item Pensad en una familia que os haya marcado positivamente. ¿Qué tenía esa familia que la hacía especial? ¿Podéis identificar formas de fecundidad que iban más allá de los hijos?

  \item ¿Cómo imagináis vuestra familia dentro de diez años? No solo en número de hijos, sino en estilo de vida: ¿será vuestra casa un lugar abierto o cerrado? ¿Qué os gustaría que la gente sintiera al cruzar vuestra puerta?

  \item La imagen de la sal y la luz: ¿en qué ámbitos concretos ---amigos, barrio, parroquia, trabajo--- creéis que vuestro matrimonio puede «dar sabor» y «dar luz»?
\end{enumerate}

\vinetacapitulo{ch10}
